\documentclass{article}
\usepackage{amsthm}
\usepackage{amsmath}
\usepackage{amssymb}

\newtheorem{theorem}{Theorem}
\newtheorem{lemma}{Lemma}
\newtheorem{proposition}{Proposition}
\newtheorem{corollary}{Corollary}
\theoremstyle{definition}
\newtheorem{definition}{Definition}
\newtheorem{example}{Example}
\theoremstyle{remark}
\newtheorem{remark}{Remark}

\begin{document}

\section{Continuity and Uniform Continuity}

\begin{definition}\label{def:continuous}
A function $f : \mathbb{R} \to \mathbb{R}$ is \emph{continuous} at a point
$a \in \mathbb{R}$ if for every $\varepsilon > 0$ there exists $\delta > 0$
such that $|f(x) - f(a)| < \varepsilon$ whenever $|x - a| < \delta$.
\end{definition}

\begin{definition}\label{def:uniform-continuous}
A function $f : \mathbb{R} \to \mathbb{R}$ is \emph{uniformly continuous}
if for every $\varepsilon > 0$ there exists $\delta > 0$ such that
$|f(x) - f(y)| < \varepsilon$ whenever $|x - y| < \delta$,
for all $x, y \in \mathbb{R}$.
\end{definition}

\begin{theorem}\label{thm:uniform-implies-continuous}
If $f$ is uniformly continuous on $\mathbb{R}$, then $f$ is continuous at
every point $a \in \mathbb{R}$.
\end{theorem}

\begin{proof}
Let $\varepsilon > 0$. Since $f$ is uniformly continuous, there exists
$\delta > 0$ such that $|f(x) - f(y)| < \varepsilon$ whenever $|x - y| < \delta$.
In particular, fixing $y = a$, we get $|f(x) - f(a)| < \varepsilon$
whenever $|x - a| < \delta$.
\end{proof}

\section{Compactness}

\begin{definition}\label{def:compact}
A subset $K \subseteq \mathbb{R}^n$ is \emph{compact} if every open cover
of $K$ has a finite subcover.
\end{definition}

\begin{theorem}\label{thm:heine-borel}
A subset $K \subseteq \mathbb{R}^n$ is compact if and only if it is
closed and bounded.
\end{theorem}

\begin{proof}
This is the Heine--Borel theorem. The proof proceeds by showing that
closed and bounded subsets of $\mathbb{R}^n$ satisfy the finite subcover
property using the completeness of $\mathbb{R}$.
\end{proof}

\begin{corollary}\label{cor:compact-continuous}
A continuous function on a compact set attains its maximum and minimum.
\end{corollary}

\begin{lemma}\label{lem:uniform-on-compact}
A continuous function $f$ on a compact set $K$ is uniformly continuous.
\end{lemma}

\section{Integration}

\begin{definition}\label{def:riemann-integral}
The \emph{Riemann integral} of $f$ over $[a, b]$ is defined as
\[
    \int_a^b f(x)\, dx = \lim_{\|\mathcal{P}\| \to 0} \sum_{i=1}^{n} f(\xi_i)\,\Delta x_i
\]
where $\mathcal{P}$ is a partition of $[a, b]$ and $\xi_i$ are sample points.
\end{definition}

\begin{theorem}\label{thm:ftc}
Let $f$ be continuous on $[a, b]$ and let $F(x) = \int_a^x f(t)\, dt$.
Then $F$ is differentiable and $F'(x) = f(x)$ for all $x \in [a, b]$.
\end{theorem}

\begin{remark}
This is the Fundamental Theorem of Calculus. It connects differentiation
and integration as inverse operations.
\end{remark}

\end{document}
